% Blueprint: Kuznetsov's theorem
% Created by Numina

\begin{document}

\section{Kuznetsov's theorem: finitely presented simple groups have solvable word problem}

We prove Kuznetsov's theorem: a finitely presented \emph{simple} group has a
decidable word problem. Fix $n \in \mathbb{N}$, a group $G$ that is simple
(hence nontrivial), and a homomorphism $\varphi : \mathrm{FreeGroup}(\mathrm{Fin}\,n)
\to G$ that is surjective and whose kernel is the normal closure of a finite set
of relators (i.e. $G$ is finitely presented through $\varphi$). Words are signed
letters $w : \mathrm{List}(\mathrm{Fin}\,n \times \mathrm{Bool})$, and
$\mathrm{FreeGroup.mk}\,w$ denotes the free-group element they spell.

The word problem predicate is $P(w) :\Leftrightarrow \varphi(\mathrm{FreeGroup.mk}\,w) = 1$,
and \emph{solvability} of the word problem means $P$ is a computable predicate
(\texttt{ComputablePred}). The strategy is Kuznetsov's r.e.-from-both-sides
argument: $P$ is recursively enumerable, its complement is recursively enumerable
(here is where simplicity is used), and a predicate that is r.e. with r.e.
complement is computable (Post's theorem, \texttt{ComputablePred.computable\_iff\_re\_compl\_re}).

\subsection{The word problem predicate}

\begin{definition}[Word problem predicate]
    \label{def:word-pred}
    \lean{LeanEval.GroupTheory.BooneHigmanSimpleProblem.wordProblemPred}
    \leanfile{LeanEval/GroupTheory/BooneHigmanSimple.lean}
    The \emph{word problem predicate} of $\varphi$ is
    \[
        P : \mathrm{List}(\mathrm{Fin}\,n \times \mathrm{Bool}) \to \mathrm{Prop},
        \qquad P(w) := \bigl(\varphi(\mathrm{FreeGroup.mk}\,w) = 1\bigr).
    \]
    Solvability of the word problem is the statement
    $\mathrm{ComputablePred}\,P$.
\end{definition}

\begin{lemma}[Word problem via the kernel]
    \label{lem:pred-iff-mem-ker}
    \lean{LeanEval.GroupTheory.BooneHigmanSimpleProblem.pred_iff_mem_ker}
    \leanfile{LeanEval/GroupTheory/BooneHigmanSimple.lean}
    \uses{def:word-pred}
    \leanok
    For every word $w$, $P(w)$ holds iff
    $\mathrm{FreeGroup.mk}\,w \in \ker\varphi$.
\end{lemma}
\begin{proof}
    \uses{def:word-pred}
    By definition of the kernel, $x \in \ker\varphi \Leftrightarrow \varphi(x) = 1$
    (\texttt{MonoidHom.mem\_ker}); apply this to $x = \mathrm{FreeGroup.mk}\,w$.
\end{proof}

\subsection{Recursive enumerability of normal-closure membership}

The engine of both sides is that membership of a free-group word in the normal
closure of a finite list of relators is recursively enumerable: a word lies in
the normal closure exactly when it can be spelled as a finite product of
conjugates of the relators, and such a finite \emph{certificate} can be searched
for by an unbounded search.

\begin{definition}[Conjugacy certificate]
    \label{def:certificate}
    \lean{LeanEval.GroupTheory.BooneHigmanSimpleProblem.evalCert}
    \leanfile{LeanEval/GroupTheory/BooneHigmanSimple/Computability.lean}
    Let $R : \mathrm{List}\,(\mathrm{List}(\mathrm{Fin}\,n \times \mathrm{Bool}))$
    be a finite list of relator words. A \emph{certificate} is a finite list of
    triples $(g_j, \varepsilon_j, k_j)$ with $g_j$ a word, $\varepsilon_j \in
    \mathrm{Bool}$ a sign, and $k_j$ an index into $R$. Its \emph{value word} is
    the concatenation
    \[
        \mathrm{eval}_R(c) := \bigl(g_1\,(R_{k_1})^{\varepsilon_1}\,g_1^{-1}\bigr)
        \cdots \bigl(g_m\,(R_{k_m})^{\varepsilon_m}\,g_m^{-1}\bigr),
    \]
    where $u^{-1}$ is the formal word inverse (\texttt{FreeGroup.invRev}) and
    $u^{+}=u$. The type of certificates is encodable
    (\texttt{Primcodable}), being built from lists, products, $\mathrm{Bool}$ and
    $\mathrm{Fin}$.
\end{definition}

\begin{lemma}[Letter equality is primitive recursive]
    \label{lem:primrec-eq-letter}
    \lean{LeanEval.GroupTheory.BooneHigmanSimpleProblem.primrec_eq_letter}
    \leanfile{LeanEval/GroupTheory/BooneHigmanSimple/Computability.lean}
    \leanok
    Equality of signed letters $a, b : \mathrm{Fin}\,n \times \mathrm{Bool}$, and
    the cancellation test $a = (b.1, \lnot\, b.2)$ deciding whether $a$ cancels
    against $b$, are primitive-recursive relations.
\end{lemma}
\begin{proof}
    The type $\mathrm{Fin}\,n \times \mathrm{Bool}$ is encodable
    (\texttt{Primcodable.prod}, \texttt{Primcodable.fin},
    \texttt{Primcodable.bool}) with decidable equality, so equality is primitive
    recursive (\texttt{Primrec.eq}). The cancellation test composes this equality
    with the primitive-recursive projections and boolean negation
    (\texttt{Primrec.fst}, \texttt{Primrec.snd}, \texttt{Primrec.not}).
\end{proof}

\begin{lemma}[The reduction step is primitive recursive]
    \label{lem:primrec-reduceStep}
    \lean{LeanEval.GroupTheory.BooneHigmanSimpleProblem.primrec_reduceStep}
    \leanfile{LeanEval/GroupTheory/BooneHigmanSimple/Computability.lean}
    \uses{lem:primrec-eq-letter}
    \leanok
    The single step of free-group reduction, which prepends a letter $a$ to an
    already-reduced word $u$ and cancels it against the head of $u$ when they are
    inverse (the function folded by \texttt{FreeGroup.reduce}), is a
    primitive-recursive function of $(a, u)$.
\end{lemma}
\begin{proof}
    \uses{lem:primrec-eq-letter}
    The step is a case analysis on whether $u$ is empty and, if not, whether its
    head cancels against $a$ (Lemma~\ref{lem:primrec-eq-letter}); in the
    cancelling case it returns the tail of $u$, otherwise $a :: u$. This is built
    from list case analysis, head/tail access and a primitive-recursive branch
    (\texttt{Primrec.list\_casesOn}, \texttt{Primrec.list\_head?},
    \texttt{Primrec.list\_tail}, \texttt{Primrec.ite}).
\end{proof}

\begin{lemma}[Reduction is primitive recursive]
    \label{lem:primrec-reduce}
    \lean{LeanEval.GroupTheory.BooneHigmanSimpleProblem.primrec_reduce}
    \leanfile{LeanEval/GroupTheory/BooneHigmanSimple/Computability.lean}
    \uses{lem:primrec-reduceStep}
    \leanok
    The map $w \mapsto \mathrm{FreeGroup.reduce}\,w$ on words is primitive
    recursive.
\end{lemma}
\begin{proof}
    \uses{lem:primrec-reduceStep}
    $\mathrm{FreeGroup.reduce}$ is the fold of the reduction step
    (Lemma~\ref{lem:primrec-reduceStep}) over the input list, so it is obtained
    from that step by primitive list recursion (\texttt{Primrec.list\_rec}) over
    the encodable type of words (\texttt{Primcodable.list}).
\end{proof}

\begin{lemma}[Reduction is computable]
    \label{lem:reduce-computable}
    \lean{LeanEval.GroupTheory.BooneHigmanSimpleProblem.reduce_computable}
    \leanfile{LeanEval/GroupTheory/BooneHigmanSimple/Computability.lean}
    \uses{lem:primrec-reduce}
    \leanok
    The map $w \mapsto \mathrm{FreeGroup.reduce}\,w$ on words is a computable
    function, and word equality after reduction is decidable.
\end{lemma}
\begin{proof}
    \uses{lem:primrec-reduce}
    Reduction is primitive recursive by Lemma~\ref{lem:primrec-reduce}, hence
    computable (\texttt{Primrec.to\_comp}). Equality of the resulting reduced
    words is decidable by \texttt{FreeGroup.instDecidableEq}.
\end{proof}

\begin{lemma}[Formal inversion is primitive recursive]
    \label{lem:primrec-invRev}
    \lean{LeanEval.GroupTheory.BooneHigmanSimpleProblem.primrec_invRev}
    \leanfile{LeanEval/GroupTheory/BooneHigmanSimple/Computability.lean}
    \leanok
    The formal word inverse $u \mapsto \mathrm{FreeGroup.invRev}\,u$, which
    reverses a word and flips the sign of each letter, is primitive recursive.
\end{lemma}
\begin{proof}
    $\mathrm{FreeGroup.invRev}\,u = \mathrm{reverse}\,(\mathrm{map}\,(\lambda
    (i,b).\,(i, \lnot b))\,u)$. Sign-flipping a letter is primitive recursive
    (\texttt{Primrec.fst}, \texttt{Primrec.snd}, \texttt{Primrec.not}), so the
    map over the list is (\texttt{Primrec.list\_map}), as is the reversal
    (\texttt{Primrec.list\_reverse}).
\end{proof}

\begin{lemma}[A conjugacy term is primitive recursive]
    \label{lem:primrec-conjTerm}
    \lean{LeanEval.GroupTheory.BooneHigmanSimpleProblem.primrec_conjTerm}
    \leanfile{LeanEval/GroupTheory/BooneHigmanSimple/Computability.lean}
    \uses{def:certificate, lem:primrec-invRev}
    \leanok
    Taking the relator list $R$ as an argument, the map sending a triple
    $(g, \varepsilon, k)$ to the word $g\,(R_k)^{\varepsilon}\,g^{-1}$ is a
    primitive-recursive function of $(R, (g,\varepsilon,k))$.
\end{lemma}
\begin{proof}
    \uses{def:certificate, lem:primrec-invRev}
    Indexing $R_k$ (with a default for out-of-range $k$) is primitive recursive
    (\texttt{Primrec.list\_getElem?}), as is choosing $R_k$ or its inverse
    according to $\varepsilon$ (Lemma~\ref{lem:primrec-invRev},
    \texttt{Primrec.cond}). The result is the concatenation of $g$, that word,
    and $\mathrm{invRev}\,g$ (Lemma~\ref{lem:primrec-invRev}), formed by
    \texttt{Primrec.list\_append}.
\end{proof}

\begin{lemma}[Certificate evaluation is primitive recursive]
    \label{lem:primrec-eval}
    \lean{LeanEval.GroupTheory.BooneHigmanSimpleProblem.primrec_eval}
    \leanfile{LeanEval/GroupTheory/BooneHigmanSimple/Computability.lean}
    \uses{def:certificate, lem:primrec-conjTerm}
    \leanok
    Taking the relator list $R$ as an argument, the map
    $(R, c) \mapsto \mathrm{eval}_R(c)$ is primitive recursive.
\end{lemma}
\begin{proof}
    \uses{def:certificate, lem:primrec-conjTerm}
    $\mathrm{eval}_R(c)$ is the flattening of the list obtained by mapping the
    conjugacy-term function (Lemma~\ref{lem:primrec-conjTerm}) over the triples of
    $c$: a primitive-recursive map (\texttt{Primrec.list\_map}) followed by
    \texttt{Primrec.list\_flatten}.
\end{proof}

\begin{lemma}[Certificate evaluation is computable]
    \label{lem:eval-computable}
    \lean{LeanEval.GroupTheory.BooneHigmanSimpleProblem.eval_computable}
    \leanfile{LeanEval/GroupTheory/BooneHigmanSimple/Computability.lean}
    \uses{def:certificate, lem:reduce-computable, lem:primrec-eval}
    Taking the relator list $R$ as an explicit argument, the relation
    \[
        \mathrm{Check}(R, w, c) :\Leftrightarrow
        \mathrm{FreeGroup.reduce}(\mathrm{eval}_R(c)) =
        \mathrm{FreeGroup.reduce}\,w
    \]
    is a computable relation in $(R, w, c)$, and for each fixed $(R, w)$ it is a
    decidable predicate in $c$. In particular this remains computable when $R$ is
    itself a computable function of the input word, e.g. $R \mathbin{+\!\!+} [w]$.
\end{lemma}
\begin{proof}
    \uses{lem:reduce-computable, lem:primrec-eval}
    By Lemma~\ref{lem:primrec-eval} the word $\mathrm{eval}_R(c)$ is a primitive
    recursive, hence computable, function of $(R, c)$. Composing with
    $\mathrm{FreeGroup.reduce}$ (Lemma~\ref{lem:reduce-computable}) on both sides
    and testing equality of reduced words (\texttt{Primrec.eq} via
    \texttt{FreeGroup.instDecidableEq}) makes $\mathrm{Check}$ a computable,
    pointwise-decidable relation. Allowing $R$ to vary computably with the input
    is composition with that computable argument (\texttt{Computable.comp}).
\end{proof}

\begin{definition}[Signed conjugate and the factor-list encoding]
    \label{def:signed-conj}
    \lean{LeanEval.GroupTheory.BooneHigmanSimpleProblem.signedConj}
    \leanfile{LeanEval/GroupTheory/BooneHigmanSimple/NormalClosure.lean}
    For $g, s$ in a group $G$ and $\varepsilon : \mathrm{Bool}$, the \emph{signed
    conjugate} is
    \[
        \mathrm{sconj}(g, \varepsilon, s) := g\,s^{\varepsilon}\,g^{-1},
        \qquad s^{\mathrm{true}} := s,\quad s^{\mathrm{false}} := s^{-1}.
    \]
    Given a set $S \subseteq G$, a \emph{factor list} over $S$ is a list
    $L : \mathrm{List}\,(G \times \mathrm{Bool} \times S)$ of factor data; its
    \emph{value} is the group product
    \[
        \mathrm{sconjProd}(L) :=
        \bigl(L.\mathrm{map}\,(\lambda (g,\varepsilon,s).\,
        \mathrm{sconj}(g,\varepsilon,s.\mathrm{val}))\bigr).\mathrm{prod}.
    \]
    The carrier of $H_S$ is the set of values of factor lists over $S$, i.e.
    $H_S := \{\, \mathrm{sconjProd}(L) : L \text{ a factor list over } S \,\}$.
    This explicit \texttt{List}-of-factor-data encoding is the carrier used in the
    three lemmas below.
\end{definition}

\begin{lemma}[Inverse of a signed conjugate]
    \label{lem:conjTerm-inv}
    \lean{LeanEval.GroupTheory.BooneHigmanSimpleProblem.signedConj_inv}
    \leanfile{LeanEval/GroupTheory/BooneHigmanSimple/NormalClosure.lean}
    \uses{def:signed-conj}
    For $g, s \in G$ and $\varepsilon : \mathrm{Bool}$,
    \[
        \mathrm{sconj}(g, \varepsilon, s)^{-1} =
        \mathrm{sconj}(g, \lnot\varepsilon, s).
    \]
\end{lemma}
\begin{proof}
    \uses{def:signed-conj}
    The inverse distributes as $(g\,s^{\varepsilon}\,g^{-1})^{-1} =
    g\,(s^{\varepsilon})^{-1}\,g^{-1}$ (\texttt{mul\_inv\_rev}, \texttt{inv\_inv}).
    A case split on $\varepsilon$ gives $(s^{\varepsilon})^{-1} = s^{\lnot\varepsilon}$
    ($\varepsilon = \mathrm{false}$ via \texttt{inv\_inv}), so the right-hand side is
    $\mathrm{sconj}(g, \lnot\varepsilon, s)$.
\end{proof}

\begin{lemma}[The carrier is closed under inverse]
    \label{lem:carrier-inv-closed}
    \lean{LeanEval.GroupTheory.BooneHigmanSimpleProblem.signedConjProd_inv_mem}
    \leanfile{LeanEval/GroupTheory/BooneHigmanSimple/NormalClosure.lean}
    \uses{def:signed-conj, lem:conjTerm-inv}
    If $x \in H_S$, i.e. $x = \mathrm{sconjProd}(L)$ for a factor list $L$ over $S$,
    then $x^{-1} \in H_S$.
\end{lemma}
\begin{proof}
    \uses{def:signed-conj, lem:conjTerm-inv}
    The inverse of a list product is the product of the reversed list of inverted
    factors (\texttt{List.prod\_inv\_reverse}). By Lemma~\ref{lem:conjTerm-inv}
    each inverted factor $\mathrm{sconj}(g,\varepsilon,s)^{-1} =
    \mathrm{sconj}(g,\lnot\varepsilon,s)$ is again a signed conjugate with the same
    $s \in S$. Hence $x^{-1} = \mathrm{sconjProd}(L')$ for the factor list $L'$
    obtained from $L$ by reversing and flipping each sign
    (\texttt{List.map\_map}, \texttt{List.map\_reverse}), so $x^{-1} \in H_S$.
\end{proof}

\begin{lemma}[Products of signed conjugates form a subgroup]
    \label{lem:prod-conjugates-subgroup}
    \lean{LeanEval.GroupTheory.BooneHigmanSimpleProblem.prodConjugatesSubgroup}
    \leanfile{LeanEval/GroupTheory/BooneHigmanSimple/NormalClosure.lean}
    \uses{def:signed-conj, lem:carrier-inv-closed}
    \leanok
    Let $S$ be a set in a group $G$. The carrier $H_S$ of
    Definition~\ref{def:signed-conj} is a subgroup of $G$.
\end{lemma}
\begin{proof}
    \uses{def:signed-conj, lem:carrier-inv-closed}
    Bundle the three subgroup fields. The identity is the value of the empty
    factor list, so $1 \in H_S$ (\texttt{one\_mem} via \texttt{List.prod\_nil}).
    For \texttt{mul\_mem}, the concatenation of two factor lists is a factor list
    whose value is the product of the two values (\texttt{List.prod\_append}),
    giving $xy \in H_S$. For \texttt{inv\_mem}, invoke
    Lemma~\ref{lem:carrier-inv-closed}.
\end{proof}

\begin{lemma}[Conjugates lie in the subgroup of products]
    \label{lem:conjugates-subset-prod}
    \lean{LeanEval.GroupTheory.BooneHigmanSimpleProblem.conjugates_subset_prod}
    \leanfile{LeanEval/GroupTheory/BooneHigmanSimple/NormalClosure.lean}
    \uses{lem:prod-conjugates-subgroup}
    \leanok
    With $H_S$ as in Lemma~\ref{lem:prod-conjugates-subgroup},
    $\mathrm{conjugatesOfSet}\,S \subseteq H_S$.
\end{lemma}
\begin{proof}
    \uses{lem:prod-conjugates-subgroup}
    An element of $\mathrm{conjugatesOfSet}\,S$ is conjugate to some $s \in S$
    (\texttt{Group.mem\_conjugatesOfSet\_iff}), i.e. of the form $g\,s\,g^{-1}$.
    This is a single sign-$+$ term, hence lies in $H_S$.
\end{proof}

\begin{lemma}[Closure elements are products of signed conjugates]
    \label{lem:closure-prod-conjugates}
    \lean{LeanEval.GroupTheory.BooneHigmanSimpleProblem.closure_prod_conjugates}
    \leanfile{LeanEval/GroupTheory/BooneHigmanSimple/NormalClosure.lean}
    \uses{lem:prod-conjugates-subgroup, lem:conjugates-subset-prod}
    \leanok
    Let $S$ be a set in a group. Every element of
    $\mathrm{normalClosure}\,S = \mathrm{closure}(\mathrm{conjugatesOfSet}\,S)$ is
    a finite product of terms of the form $g\,s^{\pm 1}\,g^{-1}$ with $s \in S$.
\end{lemma}
\begin{proof}
    \uses{lem:prod-conjugates-subgroup, lem:conjugates-subset-prod}
    By Lemma~\ref{lem:conjugates-subset-prod},
    $\mathrm{conjugatesOfSet}\,S \subseteq H_S$, and $H_S$ is a subgroup by
    Lemma~\ref{lem:prod-conjugates-subgroup}; therefore
    $\mathrm{closure}(\mathrm{conjugatesOfSet}\,S) \le H_S$
    (\texttt{Subgroup.closure\_le}). Since
    $\mathrm{normalClosure}\,S = \mathrm{closure}(\mathrm{conjugatesOfSet}\,S)$ by
    the definition of \texttt{Subgroup.normalClosure}, every element of
    $\mathrm{normalClosure}\,S$ lies in $H_S$, i.e. is such a product.
\end{proof}

\begin{lemma}[A single conjugacy term under \texttt{mk}]
    \label{lem:mk-conjTerm}
    \lean{LeanEval.GroupTheory.BooneHigmanSimpleProblem.mk_conjTerm}
    \leanfile{LeanEval/GroupTheory/BooneHigmanSimple/NormalClosure.lean}
    \uses{def:certificate}
    \leanok
    Let $\mathrm{conjTerm}_R(g, \varepsilon, k) := g \mathbin{+\!\!+}
    (R_k)^{\varepsilon} \mathbin{+\!\!+} \mathrm{FreeGroup.invRev}\,g$ be the word
    formed by Lemma~\ref{lem:primrec-conjTerm}, where $(R_k)^{+} = R_k$ and
    $(R_k)^{-} = \mathrm{FreeGroup.invRev}\,R_k$. Then
    \[
        \mathrm{FreeGroup.mk}\bigl(\mathrm{conjTerm}_R(g,\varepsilon,k)\bigr)
        = \mathrm{FreeGroup.mk}(g)\,
            (\mathrm{FreeGroup.mk}\,R_k)^{\pm 1}\,
            \mathrm{FreeGroup.mk}(g)^{-1}.
    \]
\end{lemma}
\begin{proof}
    \uses{def:certificate}
    The word $\mathrm{conjTerm}_R(g,\varepsilon,k)$ is a concatenation, and
    $\mathrm{FreeGroup.mk}$ turns concatenation into the group product:
    $\mathrm{FreeGroup.mk}(L_1 \mathbin{+\!\!+} L_2) =
    \mathrm{FreeGroup.mk}\,L_1 \cdot \mathrm{FreeGroup.mk}\,L_2$
    (\texttt{FreeGroup.mul\_mk}). The last factor is
    $\mathrm{FreeGroup.mk}(\mathrm{invRev}\,g) = \mathrm{FreeGroup.mk}(g)^{-1}$
    (\texttt{FreeGroup.mk\_invRev}), and likewise the middle factor is
    $(\mathrm{FreeGroup.mk}\,R_k)^{\pm 1}$ according to $\varepsilon$ (sign $-$
    again via \texttt{FreeGroup.mk\_invRev}). Combining gives the stated identity.
\end{proof}

\begin{lemma}[The value word as a list product of conjugacy terms]
    \label{lem:mk-eval-eq-map-prod}
    \lean{LeanEval.GroupTheory.BooneHigmanSimpleProblem.mk_eval_eq_map_prod}
    \leanfile{LeanEval/GroupTheory/BooneHigmanSimple/NormalClosure.lean}
    \uses{def:certificate}
    \leanok
    For a certificate $c$ over a relator list $R$,
    \[
        \mathrm{FreeGroup.mk}(\mathrm{eval}_R(c)) =
        \bigl(c.\mathrm{map}\,(\mathrm{FreeGroup.mk} \circ \mathrm{conjTerm}_R)\bigr).\mathrm{prod}.
    \]
\end{lemma}
\begin{proof}
    \uses{def:certificate}
    $\mathrm{eval}_R(c)$ is the flattening of $c.\mathrm{map}\,\mathrm{conjTerm}_R$.
    The map $\mathrm{FreeGroup.mk}$ is a monoid homomorphism, so it sends the
    flattened concatenation to the list product of the images of the parts
    (\texttt{FreeGroup.mul\_mk} packaged as \texttt{map\_list\_prod} over the
    flattened list, with \texttt{List.map\_map}). This is exactly
    $(c.\mathrm{map}\,(\mathrm{FreeGroup.mk} \circ \mathrm{conjTerm}_R)).\mathrm{prod}$.
\end{proof}

\begin{lemma}[The value word spells a product of conjugates]
    \label{lem:mk-eval-eq-prod}
    \lean{LeanEval.GroupTheory.BooneHigmanSimpleProblem.mk_eval_eq_prod}
    \leanfile{LeanEval/GroupTheory/BooneHigmanSimple/NormalClosure.lean}
    \uses{def:certificate, lem:mk-eval-eq-map-prod, lem:mk-conjTerm}
    \leanok
    For a certificate $c = ((g_j, \varepsilon_j, k_j))_j$ over a relator list $R$,
    \[
        \mathrm{FreeGroup.mk}(\mathrm{eval}_R(c)) =
        \prod_j \mathrm{FreeGroup.mk}(g_j)\,
            (\mathrm{FreeGroup.mk}\,R_{k_j})^{\pm 1}\,
            \mathrm{FreeGroup.mk}(g_j)^{-1},
    \]
    each factor being a signed conjugate of a relator element
    $\mathrm{FreeGroup.mk}\,R_{k_j} \in S$.
\end{lemma}
\begin{proof}
    \uses{lem:mk-eval-eq-map-prod, lem:mk-conjTerm}
    By Lemma~\ref{lem:mk-eval-eq-map-prod},
    $\mathrm{FreeGroup.mk}(\mathrm{eval}_R(c)) =
    (c.\mathrm{map}\,(\mathrm{FreeGroup.mk} \circ \mathrm{conjTerm}_R)).\mathrm{prod}$.
    Rewriting each mapped factor by Lemma~\ref{lem:mk-conjTerm}
    (\texttt{List.map\_congr\_left}) replaces
    $\mathrm{FreeGroup.mk}(\mathrm{conjTerm}_R(g_j,\varepsilon_j,k_j))$ by
    $\mathrm{FreeGroup.mk}(g_j)\,(\mathrm{FreeGroup.mk}\,R_{k_j})^{\pm 1}\,
    \mathrm{FreeGroup.mk}(g_j)^{-1}$, yielding the stated product of signed
    conjugates.
\end{proof}

\begin{lemma}[Relator elements are indexed]
    \label{lem:relator-index}
    \lean{LeanEval.GroupTheory.BooneHigmanSimpleProblem.relator_index}
    \leanfile{LeanEval/GroupTheory/BooneHigmanSimple/NormalClosure.lean}
    \leanok
    Let $S := \{\mathrm{FreeGroup.mk}\,r : r \in R\}$. For every $s \in S$ there is
    an index $k$ into $R$ with $\mathrm{FreeGroup.mk}\,R_k = s$.
\end{lemma}
\begin{proof}
    By definition of $S$, $s = \mathrm{FreeGroup.mk}\,r$ for some $r \in R$
    (\texttt{List.mem\_map}). Membership in the list $R$ gives an index $k$ with
    $R_k = r$ (\texttt{List.mem\_iff\_getElem}); then
    $\mathrm{FreeGroup.mk}\,R_k = \mathrm{FreeGroup.mk}\,r = s$.
\end{proof}

\begin{lemma}[A signed conjugate as the image of a conjugacy term]
    \label{lem:factor-eq-mk-conjTerm}
    \lean{LeanEval.GroupTheory.BooneHigmanSimpleProblem.factor_eq_mk_conjTerm}
    \leanfile{LeanEval/GroupTheory/BooneHigmanSimple/NormalClosure.lean}
    \uses{def:certificate, lem:mk-conjTerm, lem:relator-index}
    Let $S := \{\mathrm{FreeGroup.mk}\,r : r \in R\}$. For $g \in
    \mathrm{FreeGroup}(\mathrm{Fin}\,n)$, $\varepsilon : \mathrm{Bool}$ and $s \in
    S$, there is an index $k$ into $R$ with
    \[
        g\, s^{\pm 1}\, g^{-1} =
        \mathrm{FreeGroup.mk}\bigl(\mathrm{conjTerm}_R(g.\mathrm{toWord},
        \varepsilon, k)\bigr).
    \]
\end{lemma}
\begin{proof}
    \uses{lem:mk-conjTerm, lem:relator-index}
    By Lemma~\ref{lem:relator-index} choose $k$ with $\mathrm{FreeGroup.mk}\,R_k =
    s$. By Lemma~\ref{lem:mk-conjTerm},
    $\mathrm{FreeGroup.mk}(\mathrm{conjTerm}_R(g.\mathrm{toWord}, \varepsilon, k)) =
    \mathrm{FreeGroup.mk}(g.\mathrm{toWord})\,(\mathrm{FreeGroup.mk}\,R_k)^{\pm 1}\,
    \mathrm{FreeGroup.mk}(g.\mathrm{toWord})^{-1}$. Since
    $\mathrm{FreeGroup.mk}(g.\mathrm{toWord}) = g$ (\texttt{FreeGroup.mk\_toWord})
    and $\mathrm{FreeGroup.mk}\,R_k = s$, this equals $g\, s^{\pm 1}\, g^{-1}$.
\end{proof}

\begin{lemma}[Certificate from a product of conjugates]
    \label{lem:cert-of-prod}
    \lean{LeanEval.GroupTheory.BooneHigmanSimpleProblem.cert_of_prod}
    \leanfile{LeanEval/GroupTheory/BooneHigmanSimple/NormalClosure.lean}
    \uses{def:certificate, lem:factor-eq-mk-conjTerm, lem:mk-eval-eq-prod}
    Let $S := \{\mathrm{FreeGroup.mk}\,r : r \in R\}$. If $\mathrm{FreeGroup.mk}\,w$
    is a finite product of signed conjugates $g_j\, s_j^{\pm 1}\, g_j^{-1}$ with
    $s_j \in S$, then there is a certificate $c$ with
    $\mathrm{FreeGroup.mk}(\mathrm{eval}_R(c)) = \mathrm{FreeGroup.mk}\,w$.
\end{lemma}
\begin{proof}
    \uses{lem:factor-eq-mk-conjTerm, lem:mk-eval-eq-prod}
    Apply Lemma~\ref{lem:factor-eq-mk-conjTerm} to each factor $g_j\, s_j^{\pm 1}\,
    g_j^{-1}$, obtaining an index $k_j$ with that factor equal to
    $\mathrm{FreeGroup.mk}(\mathrm{conjTerm}_R(g_j.\mathrm{toWord}, \varepsilon_j,
    k_j))$. Collect the triples $(g_j.\mathrm{toWord}, \varepsilon_j, k_j)$ into a
    certificate $c$ by a \texttt{List.map} over the factors. Then by
    Lemma~\ref{lem:mk-eval-eq-prod},
    $\mathrm{FreeGroup.mk}(\mathrm{eval}_R(c)) = \prod_j
    \mathrm{FreeGroup.mk}(\mathrm{conjTerm}_R(\dots)) =
    \prod_j g_j\, s_j^{\pm 1}\, g_j^{-1} = \mathrm{FreeGroup.mk}\,w$.
\end{proof}

\begin{lemma}[A certificate witnesses normal-closure membership]
    \label{lem:mem-of-cert}
    \lean{LeanEval.GroupTheory.BooneHigmanSimpleProblem.mem_of_cert}
    \leanfile{LeanEval/GroupTheory/BooneHigmanSimple/NormalClosure.lean}
    \uses{def:certificate, lem:mk-eval-eq-prod}
    \leanok
    Let $S := \{\mathrm{FreeGroup.mk}\,r : r \in R\}$. If a certificate $c$
    satisfies $\mathrm{FreeGroup.mk}(\mathrm{eval}_R(c)) = \mathrm{FreeGroup.mk}\,w$,
    then $\mathrm{FreeGroup.mk}\,w \in \mathrm{normalClosure}\,S$.
\end{lemma}
\begin{proof}
    \uses{lem:mk-eval-eq-prod}
    By Lemma~\ref{lem:mk-eval-eq-prod}, $\mathrm{FreeGroup.mk}\,w$ is a product of
    signed conjugates of elements $\mathrm{FreeGroup.mk}\,R_{k_j} \in S$. Each such
    element lies in $\mathrm{normalClosure}\,S$
    (\texttt{Subgroup.subset\_normalClosure}), its conjugates lie there too
    (\texttt{Subgroup.conjugatesOfSet\_subset\_normalClosure}), and the subgroup
    $\mathrm{normalClosure}\,S$ is closed under product and inverse; hence the
    whole product lies in $\mathrm{normalClosure}\,S$.
\end{proof}

\begin{lemma}[Membership in a normal closure via certificates]
    \label{lem:mem-normalClosure-cert}
    \lean{LeanEval.GroupTheory.BooneHigmanSimpleProblem.mem_normalClosure_cert}
    \leanfile{LeanEval/GroupTheory/BooneHigmanSimple/NormalClosure.lean}
    \uses{def:certificate, lem:closure-prod-conjugates, lem:cert-of-prod, lem:mem-of-cert}
    \leanok
    Let $S := \{\mathrm{FreeGroup.mk}\,r : r \in R\}$ be the (finite) set of
    relator elements. For every word $w$,
    \[
        \mathrm{FreeGroup.mk}\,w \in \mathrm{normalClosure}\,S
        \iff \exists\, c,\; \mathrm{FreeGroup.mk}(\mathrm{eval}_R(c))
              = \mathrm{FreeGroup.mk}\,w .
    \]
\end{lemma}
\begin{proof}
    \uses{lem:closure-prod-conjugates, lem:cert-of-prod, lem:mem-of-cert}
    ($\Rightarrow$) By Lemma~\ref{lem:closure-prod-conjugates},
    $\mathrm{FreeGroup.mk}\,w$ is a finite product of signed conjugates
    $g_j\, s_j^{\pm 1}\, g_j^{-1}$ with $s_j \in S$; apply
    Lemma~\ref{lem:cert-of-prod} to obtain the certificate.

    ($\Leftarrow$) This is exactly Lemma~\ref{lem:mem-of-cert}.
\end{proof}

\begin{lemma}[Decision function of a computable relation]
    \label{lem:computablePred-decide}
    \lean{LeanEval.GroupTheory.BooneHigmanSimpleProblem.computablePred_decide}
    \leanfile{LeanEval/GroupTheory/BooneHigmanSimple/Computability.lean}
    \leanok
    Let $\alpha, \beta$ be encodable types and $Q : \alpha \to \beta \to
    \mathrm{Prop}$ a computable, pointwise decidable relation. Then the boolean
    function $p \mapsto \mathrm{decide}\,Q(p.1, p.2)$ on $\alpha \times \beta$ is
    computable.
\end{lemma}
\begin{proof}
    A decidable predicate is computable iff its boolean decision function is
    computable (\texttt{computablePred\_iff\_computable\_decide}). Applying this to
    the computable predicate $p \mapsto Q(p.1, p.2)$ on $\alpha \times \beta$
    gives computability of $p \mapsto \mathrm{decide}\,Q(p.1, p.2)$.
\end{proof}

\begin{lemma}[The search kernel is computable]
    \label{lem:rproj-kernel}
    \lean{LeanEval.GroupTheory.BooneHigmanSimpleProblem.rproj_kernel}
    \leanfile{LeanEval/GroupTheory/BooneHigmanSimple/Computability.lean}
    \uses{lem:computablePred-decide}
    \leanok
    With $Q$ as above and $\beta$ encodable, the option-valued function
    \[
        k(a, m) := (\mathrm{decode}\,m : \mathrm{Option}\,\beta)\ \mathbf{bind}\
            \bigl(\lambda b.\ \mathbf{if}\ Q(a,b)\ \mathbf{then}\ \mathrm{some}\,b\
            \mathbf{else}\ \mathrm{none}\bigr)
    \]
    is a computable function of $(a, m) : \alpha \times \mathbb{N}$ (i.e.
    \texttt{Computable₂}).
\end{lemma}
\begin{proof}
    \uses{lem:computablePred-decide}
    Decoding $m$ to $\mathrm{Option}\,\beta$ is computable
    (\texttt{Computable.decode}). The inner map $b \mapsto \mathbf{if}\ Q(a,b)\
    \mathbf{then}\ \mathrm{some}\,b\ \mathbf{else}\ \mathrm{none}$ is computable in
    $(a, b)$, being an \texttt{ite} on the computable boolean
    $\mathrm{decide}\,Q(a,b)$ (Lemma~\ref{lem:computablePred-decide}). Binding the
    decoded option through this map is computable (\texttt{Computable.option\_bind},
    equivalently \texttt{Computable.bind\_decode\_iff}).
\end{proof}

\begin{lemma}[The unbounded search is partial recursive]
    \label{lem:rproj-search}
    \lean{LeanEval.GroupTheory.BooneHigmanSimpleProblem.rproj_search}
    \leanfile{LeanEval/GroupTheory/BooneHigmanSimple/Computability.lean}
    \uses{lem:rproj-kernel}
    \leanok
    With $k$ as in Lemma~\ref{lem:rproj-kernel}, the partial function
    $a \mapsto \mathrm{Nat.rfindOpt}\,(\lambda m.\ k(a,m))$ is partial recursive
    (\texttt{Partrec}).
\end{lemma}
\begin{proof}
    \uses{lem:rproj-kernel}
    Immediate from \texttt{Partrec.rfindOpt} applied to the computable kernel $k$
    (Lemma~\ref{lem:rproj-kernel}).
\end{proof}

\begin{lemma}[Domain of the search equals the existential]
    \label{lem:rproj-dom}
    \lean{LeanEval.GroupTheory.BooneHigmanSimpleProblem.rproj_dom}
    \leanfile{LeanEval/GroupTheory/BooneHigmanSimple/Computability.lean}
    \uses{lem:rproj-kernel}
    \leanok
    With $k$ as in Lemma~\ref{lem:rproj-kernel}, for every $a$,
    \[
        \bigl(\mathrm{Nat.rfindOpt}\,(\lambda m.\ k(a,m))\bigr).\mathrm{Dom}
        \iff \exists b,\ Q(a,b).
    \]
\end{lemma}
\begin{proof}
    \uses{lem:rproj-kernel}
    By \texttt{Nat.rfindOpt\_dom}, the search halts on $a$ iff there exist $m$ and
    $b$ with $b \in k(a,m)$, and by definition of $k$ this means
    $\mathrm{decode}\,m = \mathrm{some}\,b$ and $Q(a,b)$. The forward direction
    yields a $b$ with $Q(a,b)$. Conversely, given $b$ with $Q(a,b)$, take
    $m := \mathrm{encode}\,b$; then $\mathrm{decode}\,m = \mathrm{some}\,b$
    (\texttt{Encodable.encodek}), so $b \in k(a,m)$ and the search halts.
\end{proof}

\begin{lemma}[Recursive enumerability of an existential over a computable relation]
    \label{lem:re-projection}
    \lean{LeanEval.GroupTheory.BooneHigmanSimpleProblem.re_projection}
    \leanfile{LeanEval/GroupTheory/BooneHigmanSimple/Computability.lean}
    Let $\beta$ be an encodable type and $Q : \alpha \to \beta \to \mathrm{Prop}$
    a relation that is computable and pointwise decidable. Then the predicate
    $a \mapsto \exists b,\; Q(a, b)$ is recursively enumerable
    (\texttt{REPred}).
\end{lemma}
\begin{proof}
    \uses{lem:rproj-search, lem:rproj-dom}
    By Lemma~\ref{lem:rproj-search} the search
    $a \mapsto \mathrm{Nat.rfindOpt}\,(\lambda m.\ k(a,m))$ is partial recursive,
    and by Lemma~\ref{lem:rproj-dom} its domain is exactly
    $\{a : \exists b,\ Q(a,b)\}$. Since \texttt{REPred} is by definition the
    predicate cut out by the domain of a partial recursive function (via the
    \texttt{Part.assert} witness), the predicate $a \mapsto \exists b,\ Q(a,b)$ is
    \texttt{REPred}; conclude with \texttt{REPred.of\_eq} (equivalently
    \texttt{Partrec.of\_eq}).
\end{proof}

\begin{lemma}[Free-group word equality is a reduce test]
    \label{lem:mk-eq-iff-reduce}
    \lean{LeanEval.GroupTheory.BooneHigmanSimpleProblem.mk_eq_iff_reduce}
    \leanfile{LeanEval/GroupTheory/BooneHigmanSimple/Computability.lean}
    \leanok
    For words $u, v$,
    \[
        \mathrm{FreeGroup.mk}\,u = \mathrm{FreeGroup.mk}\,v
        \iff \mathrm{FreeGroup.reduce}\,u = \mathrm{FreeGroup.reduce}\,v.
    \]
\end{lemma}
\begin{proof}
    ($\Leftarrow$) Words with equal reductions name the same free-group element
    (\texttt{FreeGroup.reduce.exact}). ($\Rightarrow$) Apply $\mathrm{toWord}$ to
    both sides; since $(\mathrm{FreeGroup.mk}\,u).\mathrm{toWord} =
    \mathrm{FreeGroup.reduce}\,u$ (\texttt{FreeGroup.toWord\_mk}), this yields
    $\mathrm{FreeGroup.reduce}\,u = \mathrm{FreeGroup.reduce}\,v$.
\end{proof}

\begin{lemma}[The certificate check is a computable predicate]
    \label{lem:check-computablePred}
    \lean{LeanEval.GroupTheory.BooneHigmanSimpleProblem.check_computablePred}
    \leanfile{LeanEval/GroupTheory/BooneHigmanSimple/Computability.lean}
    \uses{lem:eval-computable}
    Fix a relator list $R$. The predicate
    \[
        (w, c) \mapsto \mathrm{FreeGroup.reduce}(\mathrm{eval}_R(c))
            = \mathrm{FreeGroup.reduce}\,w
    \]
    on $\mathrm{Word} \times \mathrm{Certificate}$ is a \texttt{ComputablePred}
    (and pointwise decidable).
\end{lemma}
\begin{proof}
    \uses{lem:eval-computable}
    This is the relation $\mathrm{Check}(R, w, c)$ of
    Lemma~\ref{lem:eval-computable} with $R$ fixed, repackaged as a predicate of
    the single pair $(w, c) : \mathrm{Word} \times \mathrm{Certificate}$. It is
    computable and pointwise decidable there, hence a \texttt{ComputablePred}.
\end{proof}

\begin{lemma}[Normal-closure membership is r.e.]
    \label{lem:re-mem-normalClosure}
    \lean{LeanEval.GroupTheory.BooneHigmanSimpleProblem.re_mem_normalClosure}
    \leanfile{LeanEval/GroupTheory/BooneHigmanSimple.lean}
    \uses{lem:mem-normalClosure-cert, lem:mk-eq-iff-reduce, lem:check-computablePred, lem:re-projection}
    Fix a relator list $R$ and $S = \{\mathrm{FreeGroup.mk}\,r : r \in R\}$. The
    predicate $w \mapsto \mathrm{FreeGroup.mk}\,w \in \mathrm{normalClosure}\,S$
    is recursively enumerable.
\end{lemma}
\begin{proof}
    \uses{lem:mem-normalClosure-cert, lem:mk-eq-iff-reduce, lem:check-computablePred, lem:re-projection}
    By Lemma~\ref{lem:mem-normalClosure-cert} the predicate is equivalent to
    $w \mapsto \exists c,\; \mathrm{FreeGroup.mk}(\mathrm{eval}_R(c)) =
    \mathrm{FreeGroup.mk}\,w$. By the bridge Lemma~\ref{lem:mk-eq-iff-reduce} each
    such equality is the reduce test
    $\mathrm{FreeGroup.reduce}(\mathrm{eval}_R(c)) = \mathrm{FreeGroup.reduce}\,w$,
    so the predicate is equivalent to $w \mapsto \exists c,\; \mathrm{Check}(R, w, c)$.
    The relation $\mathrm{Check}$ (with $R$ fixed) is a computable, pointwise
    decidable predicate of $(w, c)$ by Lemma~\ref{lem:check-computablePred}; apply
    Lemma~\ref{lem:re-projection} to the existential over $c$ and transport along
    the equivalence (\texttt{REPred.of\_eq}).
\end{proof}

\subsection{The positive side: the word problem is r.e.}

\begin{lemma}[A finite set of free-group elements is spelled by a word list]
    \label{lem:finset-to-word-list}
    \lean{LeanEval.GroupTheory.BooneHigmanSimpleProblem.finset_to_word_list}
    \leanfile{LeanEval/GroupTheory/BooneHigmanSimple/NormalClosure.lean}
    \leanok
    For every finite set $T \subseteq \mathrm{FreeGroup}(\mathrm{Fin}\,n)$ there
    is a finite list of words $R$ with
    $\{\mathrm{FreeGroup.mk}\,r : r \in R\} = T$.
\end{lemma}
\begin{proof}
    Take $R := (T.\mathrm{toList}).\mathrm{map}\,(\lambda t.\, t.\mathrm{toWord})$.
    Each element $t \in T$ equals $\mathrm{FreeGroup.mk}$ of its reduced word
    (\texttt{FreeGroup.mk\_toWord}), and $T.\mathrm{toList}$ enumerates exactly
    $T$ (\texttt{Finset.mem\_toList}), so mapping recovers $T$ as the set of
    $\mathrm{FreeGroup.mk}\,r$ (\texttt{List.mem\_map}).
\end{proof}

\begin{lemma}[The kernel is a normal closure of relators]
    \label{lem:ker-eq-normalClosure}
    \lean{LeanEval.GroupTheory.BooneHigmanSimpleProblem.ker_eq_normalClosure}
    \leanfile{LeanEval/GroupTheory/BooneHigmanSimple.lean}
    \uses{lem:finset-to-word-list}
    \leanok
    There is a finite relator list $R$ such that
    $\ker\varphi = \mathrm{normalClosure}\,\{\mathrm{FreeGroup.mk}\,r : r \in R\}$.
\end{lemma}
\begin{proof}
    \uses{lem:finset-to-word-list}
    By hypothesis $\ker\varphi$ is a finitely generated normal closure
    (\texttt{Subgroup.IsNormalClosureFG}): there is a finite set
    $T \subseteq \mathrm{FreeGroup}(\mathrm{Fin}\,n)$ with
    $\mathrm{normalClosure}\,T = \ker\varphi$. By
    Lemma~\ref{lem:finset-to-word-list} there is a finite list $R$ with
    $\{\mathrm{FreeGroup.mk}\,r : r \in R\} = T$, whence the claim.
\end{proof}

\begin{lemma}[Positive side is r.e.]
    \label{lem:re-positive}
    \lean{LeanEval.GroupTheory.BooneHigmanSimpleProblem.re_positive}
    \leanfile{LeanEval/GroupTheory/BooneHigmanSimple.lean}
    \uses{def:word-pred, lem:pred-iff-mem-ker, lem:ker-eq-normalClosure, lem:re-mem-normalClosure}
    \leanok
    The word problem predicate $P$ is recursively enumerable.
\end{lemma}
\begin{proof}
    \uses{lem:pred-iff-mem-ker, lem:ker-eq-normalClosure, lem:re-mem-normalClosure}
    By Lemma~\ref{lem:pred-iff-mem-ker}, $P(w) \Leftrightarrow
    \mathrm{FreeGroup.mk}\,w \in \ker\varphi$, and by
    Lemma~\ref{lem:ker-eq-normalClosure} the kernel is the normal closure of the
    finite relator set $S$. So $P$ is pointwise equal to the predicate of
    Lemma~\ref{lem:re-mem-normalClosure}, which is r.e.; conclude with
    \texttt{REPred.of\_eq}.
\end{proof}

\subsection{The negative side: simplicity collapses the group}

\begin{lemma}[Normal closure of a nontrivial element in a simple group]
    \label{lem:normalClosure-singleton-top}
    \lean{LeanEval.GroupTheory.BooneHigmanSimpleProblem.normalClosure_singleton_top}
    \leanfile{LeanEval/GroupTheory/BooneHigmanSimple/NormalClosure.lean}
    \leanok
    If $g \in G$ and $g \neq 1$, then $\mathrm{normalClosure}\,\{g\} = \top$.
\end{lemma}
\begin{proof}
    The normal closure of any set is a normal subgroup
    (\texttt{Subgroup.normalClosure\_normal}). Since $G$ is simple, this normal
    subgroup is $\bot$ or $\top$ (\texttt{Subgroup.Normal.eq\_bot\_or\_eq\_top}).
    It is not $\bot$: $\mathrm{normalClosure}\,\{g\} = \bot$ would force
    $\{g\} \subseteq \{1\}$ (\texttt{Subgroup.normalClosure\_eq\_bot\_iff}), i.e.
    $g = 1$, contrary to hypothesis. Hence it is $\top$.
\end{proof}

\begin{lemma}[The kernel is not everything]
    \label{lem:ker-ne-top}
    \lean{LeanEval.GroupTheory.BooneHigmanSimpleProblem.ker_ne_top}
    \leanfile{LeanEval/GroupTheory/BooneHigmanSimple.lean}
    \uses{lem:ker-eq-normalClosure}
    \leanok
    $\ker\varphi \neq \top$.
\end{lemma}
\begin{proof}
    \uses{lem:ker-eq-normalClosure}
    If $\ker\varphi = \top$ then every element maps to $1$ under $\varphi$
    (\texttt{MonoidHom.ker\_eq\_top\_iff}); since $\varphi$ is surjective, $G$
    would be trivial, contradicting that $G$ is simple and therefore nontrivial
    (\texttt{IsSimpleGroup} extends \texttt{Nontrivial}).
\end{proof}

\begin{lemma}[A subgroup is everything iff it contains all generators]
    \label{lem:top-iff-generators}
    \lean{LeanEval.GroupTheory.BooneHigmanSimpleProblem.top_iff_generators}
    \leanfile{LeanEval/GroupTheory/BooneHigmanSimple/NormalClosure.lean}
    \leanok
    For a subgroup $N \le \mathrm{FreeGroup}(\mathrm{Fin}\,n)$,
    \[
        N = \top \iff \forall i : \mathrm{Fin}\,n,\; \mathrm{FreeGroup.of}\,i \in N .
    \]
\end{lemma}
\begin{proof}
    The generators $\mathrm{FreeGroup.of}\,i$ generate the free group: the
    subgroup closure of their range is $\top$
    (\texttt{FreeGroup.closure\_range\_of}). If $N$ contains all of them then
    $N \ge \mathrm{closure}(\mathrm{range}\,\mathrm{FreeGroup.of}) = \top$, so
    $N = \top$; the converse is immediate.
\end{proof}

\begin{lemma}[Surjective image criterion for being everything]
    \label{lem:top-iff-map-top}
    \lean{LeanEval.GroupTheory.BooneHigmanSimpleProblem.top_iff_map_top}
    \leanfile{LeanEval/GroupTheory/BooneHigmanSimple/NormalClosure.lean}
    \leanok
    Let $\varphi : G \to H$ be a surjective homomorphism and $N \le G$ a subgroup
    with $\ker\varphi \le N$. Then
    $N = \top \iff \mathrm{map}\,\varphi\,N = \top$.
\end{lemma}
\begin{proof}
    Mapping is monotone and $\mathrm{map}\,\varphi\,\top = \top$ by surjectivity
    (\texttt{Subgroup.map\_top\_of\_surjective}), which gives the forward
    direction. Conversely, since $\ker\varphi \le N$,
    $\mathrm{comap}\,\varphi\,(\mathrm{map}\,\varphi\,N) = N$
    (\texttt{Subgroup.comap\_map\_eq\_self}); so if $\mathrm{map}\,\varphi\,N =
    \top$ then $N = \mathrm{comap}\,\varphi\,\top = \top$.
\end{proof}

\begin{lemma}[Adjoining a subset of the identity preserves the normal closure]
    \label{lem:normalClosure-insert-subset-one}
    \lean{LeanEval.GroupTheory.BooneHigmanSimpleProblem.normalClosure_union_subset_one}
    \leanfile{LeanEval/GroupTheory/BooneHigmanSimple/NormalClosure.lean}
    Let $A, T$ be sets in a group with $T \subseteq \{1\}$. Then
    $\mathrm{normalClosure}\,(A \cup T) = \mathrm{normalClosure}\,A$.
\end{lemma}
\begin{proof}
    By antisymmetry. For $\ge$: $A \subseteq A \cup T$ gives
    $\mathrm{normalClosure}\,A \le \mathrm{normalClosure}\,(A \cup T)$ by
    monotonicity (\texttt{Subgroup.normalClosure\_mono}). For $\le$: since
    $\mathrm{normalClosure}\,A$ is normal and minimal among normal subgroups
    containing $A$ (\texttt{Subgroup.normalClosure\_subset\_iff}), it suffices that
    $A \cup T \subseteq \mathrm{normalClosure}\,A$; indeed $A \subseteq
    \mathrm{normalClosure}\,A$ (\texttt{Subgroup.subset\_normalClosure}) and
    $T \subseteq \{1\} \subseteq \mathrm{normalClosure}\,A$ as the latter contains
    $1$ (\texttt{one\_mem}). Hence both closures coincide.
\end{proof}

\begin{lemma}[Image of the extended normal closure]
    \label{lem:map-normalClosure-insert}
    \lean{LeanEval.GroupTheory.BooneHigmanSimpleProblem.map_normalClosure_insert}
    \leanfile{LeanEval/GroupTheory/BooneHigmanSimple/NormalClosure.lean}
    \uses{lem:normalClosure-insert-subset-one}
    \leanok
    Let $\varphi : G \to H$ be a surjective homomorphism, $S$ a set with
    $\mathrm{normalClosure}\,S = \ker\varphi$, and $x \in G$. Then
    \[
        \mathrm{map}\,\varphi\,\bigl(\mathrm{normalClosure}\,(S \cup \{x\})\bigr)
        = \mathrm{normalClosure}\,\{\varphi(x)\}.
    \]
\end{lemma}
\begin{proof}
    \uses{lem:normalClosure-insert-subset-one}
    Since $\varphi$ is surjective, the image of a normal closure is the normal
    closure of the image (\texttt{Subgroup.map\_normalClosure}):
    $\mathrm{map}\,\varphi\,(\mathrm{normalClosure}\,(S \cup \{x\})) =
    \mathrm{normalClosure}\,(\varphi(S) \cup \{\varphi(x)\})$. Every $s \in S$ lies
    in $\mathrm{normalClosure}\,S = \ker\varphi$, so $\varphi(s) = 1$
    (\texttt{MonoidHom.mem\_ker}); thus $\varphi(S) \subseteq \{1\}$. Apply
    Lemma~\ref{lem:normalClosure-insert-subset-one} with $A = \{\varphi(x)\}$ and
    $T = \varphi(S)$ to conclude
    $\mathrm{normalClosure}\,(\varphi(S) \cup \{\varphi(x)\}) =
    \mathrm{normalClosure}\,\{\varphi(x)\}$.
\end{proof}

\begin{lemma}[Collapse criterion]
    \label{lem:collapse}
    \lean{LeanEval.GroupTheory.BooneHigmanSimpleProblem.collapse}
    \leanfile{LeanEval/GroupTheory/BooneHigmanSimple.lean}
    \uses{lem:pred-iff-mem-ker, lem:ker-eq-normalClosure, lem:normalClosure-singleton-top, lem:ker-ne-top, lem:top-iff-map-top, lem:map-normalClosure-insert}
    Let $S = \{\mathrm{FreeGroup.mk}\,r : r \in R\}$ be the relator set with
    $\mathrm{normalClosure}\,S = \ker\varphi$. For every word $w$,
    \[
        \varphi(\mathrm{FreeGroup.mk}\,w) \neq 1
        \iff \mathrm{normalClosure}\,(S \cup \{\mathrm{FreeGroup.mk}\,w\}) = \top .
    \]
\end{lemma}
\begin{proof}
    \uses{lem:pred-iff-mem-ker, lem:ker-eq-normalClosure, lem:normalClosure-singleton-top, lem:ker-ne-top, lem:top-iff-map-top, lem:map-normalClosure-insert}
    Write $N := \mathrm{normalClosure}\,(S \cup \{\mathrm{FreeGroup.mk}\,w\})$;
    then $N \supseteq \mathrm{normalClosure}\,S = \ker\varphi$. By
    Lemma~\ref{lem:map-normalClosure-insert},
    $\mathrm{map}\,\varphi\,N =
    \mathrm{normalClosure}\,\{\varphi(\mathrm{FreeGroup.mk}\,w)\}$.

    ($\Rightarrow$) If $\varphi(\mathrm{FreeGroup.mk}\,w) \neq 1$ then
    $\mathrm{map}\,\varphi\,N = \top$ by
    Lemma~\ref{lem:normalClosure-singleton-top}. Since $\ker\varphi \le N$,
    Lemma~\ref{lem:top-iff-map-top} gives $N = \top$.

    ($\Leftarrow$) If instead $\varphi(\mathrm{FreeGroup.mk}\,w) = 1$, then
    $\mathrm{FreeGroup.mk}\,w \in \ker\varphi = \mathrm{normalClosure}\,S$, so
    adding it changes nothing and $N = \mathrm{normalClosure}\,S = \ker\varphi
    \neq \top$ by Lemma~\ref{lem:ker-ne-top}. Contrapositively, $N = \top$ forces
    $\varphi(\mathrm{FreeGroup.mk}\,w) \neq 1$.
\end{proof}

\begin{lemma}[Recursive enumerability is closed under conjunction]
    \label{lem:re-and}
    \lean{LeanEval.GroupTheory.BooneHigmanSimpleProblem.re_and}
    \leanfile{LeanEval/GroupTheory/BooneHigmanSimple/Computability.lean}
    \leanok
    If $p$ and $q$ are recursively enumerable predicates on the same encodable
    type, then $a \mapsto p(a) \wedge q(a)$ is recursively enumerable.
\end{lemma}
\begin{proof}
    A predicate is r.e. iff it is the domain of a partial recursive function.
    Running the two partial functions in sequence and returning once both halt
    (\texttt{Partrec.bind} on the \texttt{Part.assert} witnesses) yields a partial
    recursive function whose domain is $\{a : p(a) \wedge q(a)\}$; conclude with
    \texttt{REPred.of\_eq}.
\end{proof}

\begin{lemma}[Finite conjunction over generators is r.e.]
    \label{lem:re-forall-fin}
    \lean{LeanEval.GroupTheory.BooneHigmanSimpleProblem.re_forall_fin}
    \leanfile{LeanEval/GroupTheory/BooneHigmanSimple/Computability.lean}
    \uses{lem:re-and}
    \leanok
    If for each $i : \mathrm{Fin}\,n$ the predicate $p_i$ is recursively
    enumerable, then $a \mapsto \forall i : \mathrm{Fin}\,n,\; p_i(a)$ is
    recursively enumerable.
\end{lemma}
\begin{proof}
    \uses{lem:re-and}
    Induct on $n$. For $n = 0$ the universal statement is vacuously true, a
    decidable predicate, which is r.e. because every computable (in particular
    decidable) predicate is recursively enumerable (\texttt{ComputablePred.to\_re}).
    For the inductive step, a universal over
    $\mathrm{Fin}(n+1)$ is the conjunction of $p_0$ with a universal over
    $\mathrm{Fin}\,n$ (\texttt{Fin.forall\_fin\_succ}); apply
    Lemma~\ref{lem:re-and} and the inductive hypothesis, then
    \texttt{REPred.of\_eq}.
\end{proof}

\begin{lemma}[Relator set of an extended list]
    \label{lem:relatorSet-append}
    \lean{LeanEval.GroupTheory.BooneHigmanSimpleProblem.relatorSet_append}
    \leanfile{LeanEval/GroupTheory/BooneHigmanSimple/NormalClosure.lean}
    \leanok
    For a relator list $R$ and a word $w$, writing
    $S_R := \{\mathrm{FreeGroup.mk}\,r : r \in R\}$,
    \[
        S_{R \mathbin{+\!\!+} [w]} = S_R \cup \{\mathrm{FreeGroup.mk}\,w\}.
    \]
\end{lemma}
\begin{proof}
    By set extensionality. An element lies in $S_{R \mathbin{+\!\!+} [w]}$ iff it is
    $\mathrm{FreeGroup.mk}\,r$ for some $r \in R \mathbin{+\!\!+} [w]$, and
    $r \in R \mathbin{+\!\!+} [w] \iff r \in R \vee r = w$
    (\texttt{List.mem\_append}, \texttt{List.mem\_singleton}). Splitting the
    disjunction identifies this with membership in
    $S_R \cup \{\mathrm{FreeGroup.mk}\,w\}$.
\end{proof}

\begin{lemma}[A generator as a one-letter word]
    \label{lem:of-eq-mk-singleton}
    \lean{LeanEval.GroupTheory.BooneHigmanSimpleProblem.of_eq_mk_singleton}
    \leanfile{LeanEval/GroupTheory/BooneHigmanSimple/NormalClosure.lean}
    \leanok
    For $i : \mathrm{Fin}\,n$,
    $\mathrm{FreeGroup.of}\,i = \mathrm{FreeGroup.mk}\,[(i, \mathrm{true})]$.
\end{lemma}
\begin{proof}
    This is the definition of $\mathrm{FreeGroup.of}$ as the class of the
    one-letter word $[(i, \mathrm{true})]$ (\texttt{FreeGroup.of}).
\end{proof}

\begin{lemma}[Generator membership in the extended closure is r.e.]
    \label{lem:re-generator-mem}
    \lean{LeanEval.GroupTheory.BooneHigmanSimpleProblem.re_generator_mem}
    \leanfile{LeanEval/GroupTheory/BooneHigmanSimple.lean}
    \uses{lem:relatorSet-append, lem:of-eq-mk-singleton, lem:mem-normalClosure-cert, lem:eval-computable, lem:re-projection}
    For each fixed $i : \mathrm{Fin}\,n$, the predicate
    \[
        w \mapsto \mathrm{FreeGroup.of}\,i \in
            \mathrm{normalClosure}\,(S \cup \{\mathrm{FreeGroup.mk}\,w\})
    \]
    is recursively enumerable.
\end{lemma}
\begin{proof}
    \uses{lem:relatorSet-append, lem:of-eq-mk-singleton, lem:mem-normalClosure-cert, lem:eval-computable, lem:re-projection}
    By Lemma~\ref{lem:relatorSet-append} the extended relator set is
    $S \cup \{\mathrm{FreeGroup.mk}\,w\} = S_{R \mathbin{+\!\!+} [w]}$, so the normal
    closure in question is generated by the list $R \mathbin{+\!\!+} [w]$, which
    depends computably on $w$ (appending is a computable function of $w$). Writing
    the fixed generator as the one-letter word
    $\mathrm{FreeGroup.of}\,i = \mathrm{FreeGroup.mk}\,[(i,\mathrm{true})]$
    (Lemma~\ref{lem:of-eq-mk-singleton}) and applying
    Lemma~\ref{lem:mem-normalClosure-cert} to the list $R \mathbin{+\!\!+} [w]$,
    membership is equivalent to the existence of a certificate $c$ with
    $\mathrm{Check}(R \mathbin{+\!\!+} [w],\, [(i,\mathrm{true})],\, c)$. Because
    $R \mathbin{+\!\!+} [w]$ varies computably with $w$, this is a computable,
    pointwise decidable relation in $(w, c)$ by Lemma~\ref{lem:eval-computable}
    (its $R$-varies-computably form). Apply Lemma~\ref{lem:re-projection} to the
    existential over $c$ and conclude with \texttt{REPred.of\_eq}.
\end{proof}

\begin{lemma}[Negation as a universal over generators]
    \label{lem:neg-iff-forall-gen}
    \lean{LeanEval.GroupTheory.BooneHigmanSimpleProblem.neg_iff_forall_gen}
    \leanfile{LeanEval/GroupTheory/BooneHigmanSimple.lean}
    \uses{lem:collapse, lem:top-iff-generators}
    Let $S = \{\mathrm{FreeGroup.mk}\,r : r \in R\}$ be the relator set with
    $\mathrm{normalClosure}\,S = \ker\varphi$. For every word $w$,
    \[
        \varphi(\mathrm{FreeGroup.mk}\,w) \neq 1
        \iff \forall i : \mathrm{Fin}\,n,\ \mathrm{FreeGroup.of}\,i \in
            \mathrm{normalClosure}\,(S \cup \{\mathrm{FreeGroup.mk}\,w\}).
    \]
\end{lemma}
\begin{proof}
    \uses{lem:collapse, lem:top-iff-generators}
    By Lemma~\ref{lem:collapse}, $\varphi(\mathrm{FreeGroup.mk}\,w) \neq 1$ holds
    iff $\mathrm{normalClosure}\,(S \cup \{\mathrm{FreeGroup.mk}\,w\}) = \top$, and
    by Lemma~\ref{lem:top-iff-generators} the latter holds iff every generator
    $\mathrm{FreeGroup.of}\,i$ lies in that normal closure.
\end{proof}

\begin{lemma}[Negative side is r.e.]
    \label{lem:re-negative}
    \lean{LeanEval.GroupTheory.BooneHigmanSimpleProblem.re_negative}
    \leanfile{LeanEval/GroupTheory/BooneHigmanSimple.lean}
    \uses{def:word-pred, lem:neg-iff-forall-gen, lem:collapse, lem:top-iff-generators, lem:re-forall-fin, lem:re-generator-mem, lem:ker-eq-normalClosure}
    \leanok
    The complement predicate $w \mapsto \neg P(w)$, i.e.
    $\varphi(\mathrm{FreeGroup.mk}\,w) \neq 1$, is recursively enumerable.
\end{lemma}
\begin{proof}
    \uses{lem:neg-iff-forall-gen, lem:re-forall-fin, lem:re-generator-mem, lem:ker-eq-normalClosure}
    Unfolding $P$, the complement is $w \mapsto
    \varphi(\mathrm{FreeGroup.mk}\,w) \neq 1$. Fix the finite relator list $R$ with
    $S = \{\mathrm{FreeGroup.mk}\,r : r \in R\}$ and $\mathrm{normalClosure}\,S =
    \ker\varphi$ (Lemma~\ref{lem:ker-eq-normalClosure}). By
    Lemma~\ref{lem:neg-iff-forall-gen} this complement is pointwise equal to
    \[
        w \mapsto \forall i : \mathrm{Fin}\,n,\ \mathrm{FreeGroup.of}\,i \in
            \mathrm{normalClosure}\,(S \cup \{\mathrm{FreeGroup.mk}\,w\}).
    \]
    For each fixed $i$ the inner predicate is r.e. in $w$ by
    Lemma~\ref{lem:re-generator-mem}, so the finite conjunction over
    $\mathrm{Fin}\,n$ is r.e. by Lemma~\ref{lem:re-forall-fin}; transport along the
    equivalence with \texttt{REPred.of\_eq}.
\end{proof}

\subsection{Main result}

\begin{lemma}[Post's theorem packaged for the word problem]
    \label{lem:post-re-compl}
    \lean{LeanEval.GroupTheory.BooneHigmanSimpleProblem.post_re_compl}
    \leanfile{LeanEval/GroupTheory/BooneHigmanSimple.lean}
    \uses{def:word-pred}
    \leanok
    Let $P$ be the word problem predicate, equipped with a classical decidability
    instance. If $P$ is recursively enumerable and the complement
    $w \mapsto \neg P(w)$ is recursively enumerable, then
    $\mathrm{ComputablePred}\,P$.
\end{lemma}
\begin{proof}
    \uses{def:word-pred}
    Choose a $\mathrm{DecidablePred}$ instance for $P$ classically
    (\texttt{Classical.dec}). Post's theorem
    (\texttt{ComputablePred.computable\_iff\_re\_compl\_re}) states, for a
    predicate with a decidability instance, that $\mathrm{ComputablePred}\,P \iff
    \mathrm{REPred}\,P \wedge \mathrm{REPred}\,(\lambda w.\ \neg P(w))$. The
    right-hand conjunction is exactly the hypothesis, so the backward direction of
    the iff gives $\mathrm{ComputablePred}\,P$.
\end{proof}

\begin{theorem}[Kuznetsov / Boone--Higman for simple groups]
    \label{thm:boone-higman-simple}
    \lean{LeanEval.GroupTheory.BooneHigmanSimpleProblem.boone_higman_simple}
    \leanfile{LeanEval/GroupTheory/BooneHigmanSimple.lean}
    \uses{def:word-pred, lem:post-re-compl, lem:re-positive, lem:re-negative}
    \leanok
    The word problem predicate $P$ of a finitely presented simple group is
    computable; that is, $\mathrm{ComputablePred}\,P$, the solvability of the word
    problem of $\varphi$.
\end{theorem}
\begin{proof}
    \uses{lem:post-re-compl, lem:re-positive, lem:re-negative}
    Apply Lemma~\ref{lem:post-re-compl}. Its two hypotheses are that $P$ is r.e.
    (Lemma~\ref{lem:re-positive}) and that the complement $w \mapsto \neg P(w)$ is
    r.e. (Lemma~\ref{lem:re-negative}, whose statement is already in the
    complement shape $\mathrm{REPred}\,(\lambda w.\ \neg P(w))$ that Post's theorem
    expects). Hence $\mathrm{ComputablePred}\,P$.
\end{proof}

\end{document}
